\noindent\textbf{KETENTUAN KHUSUS — OUTLINE APLIKASI STARTUP (Kode 158):}
\begin{itemize}
    \item Dikerjakan secara KELOMPOK, MAKSIMAL 2 orang (bukan individu).
    \item Aplikasi yang dibuat WAJIB tersedia dan terdaftar di Google Play Store.
    \item WAJIB membuat dan mempresentasikan PITCHDECK saat sidang berlangsung.\\(Pitchdeck: presentasi singkat gambaran umum rencana bisnis startup.)
    \item Tinjauan StartUp di sub-bab 3.1 WAJIB berisi 9 komponen BMC lengkap.
\end{itemize}

\vspace{0.5cm}


\chapter{ANALISIS DAN PERANCANGAN APLIKASI}
\section{Tinjauan StartUp}


Tuliskan di sini \textbf{Pedoman Penulisan:} Berisikan 9 poin BMC: Key Partnership, Key Activities, Key Resources, Value Proposition, Customer Relationship, Channel, Customer Segments, Cost Structure, dan Revenue Stream.


\section{Analisis Kebutuhan}
\subsection{Analisis Kebutuhan Aplikasi}


Tuliskan di sini \textbf{Pedoman Penulisan:} Analisa dari sisi pengguna. Contoh: Pengguna dapat memilih produk, menjual, dan memesan.


\subsection{Rancangan Diagram Use Case}


Tuliskan di sini \textbf{Pedoman Penulisan:} Rancangan Use Case dari aplikasi.


\subsection{Rancangan Diagram Aktivitas}


Tuliskan di sini \textbf{Pedoman Penulisan:} Rancangan Diagram Aktivitas.


\subsection{Rancangan User Interface}


Tuliskan di sini \textbf{Pedoman Penulisan:} Rancangan antarmuka berisi MockUp dari aplikasi.


\subsection{Rancangan Database}


Tuliskan di sini \textbf{Pedoman Penulisan:} Rancangan database berisi ERD dan Logical Record Structure (LRS).


\section{Implementasi}
\subsection{Proses Implementasi}
\subsection{Hasil Implementasi}
\section{Spesifikasi Aplikasi}
\subsection{Hardware}
\subsection{Software}
\section{Pengujian Aplikasi}
\section{Uraian Tugas}


Tuliskan di sini \textbf{Pedoman Penulisan:} Mengingat StartUp maksimal dikerjakan 2 mahasiswa, uraikan pembagian tugas masing-masing.


\subsection{Project Manager / Sistem Analis}
\subsection{Database Administrator / Programmer}
\subsection{Pengujian Sistem}
